\section{Data Description and Collection}

This study integrates \textbf{three} independent, publicly available sources,
each describing Israeli high schools from a different angle. \textbf{(1)~Bagrut
exam records} (2013--2016), released under the Israeli Freedom of Information Law
and mirrored on Kaggle, give per-subject results for each school and year.
\textbf{(2)~The Central Bureau of Statistics (CBS) socioeconomic index}
rates every locality on a 1--10 cluster capturing municipal wealth, education,
and employment. \textbf{(3)~The Ministry of Education's school budget and
institutional report} provides per-institution funding, staffing, class size,
sector, and supervision type. Source~1 supplies the \emph{outcomes}, source~2 the
\emph{municipal} environment, and source~3 lets us look \emph{inside} the school
itself. The unit of analysis is the \emph{school-year}, not the student.

\textbf{Structure and storage.} The Bagrut file is a UTF-8 CSV. The CBS and
Ministry files are Excel workbooks with awkward layouts (a header offset, Hebrew
column names, and a malformed style block requiring a lenient reader). Each
source is ingested by a dedicated loader and cached as a clean CSV
(Table~\ref{tab:datasets}).

\begin{table}[ht]
\centering
\small
\renewcommand{\arraystretch}{1.35}
\caption{The three integrated data sources: one supplies the outcomes, one the
municipal context, and one the school's own institutional profile. This split is
what lets the study separate ``where a school is'' from ``what a school does''.}
\label{tab:datasets}
\begin{tabular}{@{}>{\raggedright\arraybackslash}p{2.9cm}>{\raggedright\arraybackslash}p{2.5cm}rr>{\raggedright\arraybackslash}p{2.9cm}l@{}}
\toprule
\textbf{Dataset} & \textbf{Source} & \textbf{Rows} & \textbf{Cols} & \textbf{Grain (key)} & \textbf{Coverage} \\
\midrule
Bagrut exam records & Israeli FOI (Kaggle) & 69,638 & 9 & subject cell (1,063 schools, 315 cities) & 2013--2016 \\
CBS socioeconomic index & Israel CBS & 1,208 & 7 & locality (cluster 1--10) & 2015 index \\
Ministry budget \& institution report & Israel Min.\ of Education & 4,718 & 26 & school (\texttt{semel}) & 2014/15 snapshot \\
\bottomrule
\end{tabular}
\end{table}

\textbf{Linking the sources.} The tables are joined on \emph{two different keys},
because \texttt{semel} is not shared consistently: in the Bagrut file it is a
\emph{school} code, in the CBS file a \emph{locality} code, giving zero overlap.
We therefore (i)~join Bagrut to CBS on a \emph{normalised Hebrew locality name}
through a four-stage cascade
(exact~$\rightarrow$~structural~$\rightarrow$~hand-verified crosswalk~$\rightarrow$~conservative
fuzzy match), integrating \textbf{99.44\%} of records, and (ii)~join Bagrut to
the Ministry report directly on \texttt{semel} (a shared \emph{school} code),
integrating \textbf{99.68\%}. A left join can multiply rows when the right-hand
key repeats, so each lookup table is de-duplicated on its join key first. This
matters because the normalised locality name is not unique in the CBS file. Row
preservation was verified rather than assumed: the merged table has exactly
\textbf{69,638} rows, matching the Bagrut input, and \textbf{98.7\%} of records
carry a complete profile from all three sources.

\textbf{Biases and challenges.} Three were material. The key incompatibility
above is resolved by the name cascade rather than a fragile single-pass match.
The Ministry report is a single-year (2014/15) snapshot broadcast across all four
Bagrut years, a documented limitation that assumes school resourcing is
reasonably stable over 2013--2016. Finally, the unmatched 0.56\% of records
(boarding and regional schools with no single home municipality) are
\emph{flagged, not forced}, avoiding an invented socioeconomic value.
